\paragraph{C.3q (Proved under the stated geometric hypotheses): lifting abelian subgroups of the local Picard functor in characteristic $p$.}
Let $k$ be algebraically closed of characteristic
$p>0$. Let $R$ be the strict henselization at a
$k$-rational point of a normal local scheme
essentially of finite type over $k$. Suppose
$N=\dim R\ge4$ and that
\[
 X=\operatorname{Spec}R\setminus\{\mathfrak m\},
 \qquad Y=X\cap V(f),\qquad0\ne f\in\mathfrak m,
\]
are regular. Put
\[
 G=\mathbf{Pic}^{\mathrm{loc}}_{R/k},\qquad
 H=\prod_a\mathbf{Pic}^{\mathrm{loc}}_{C_a/k},
 \qquad r:G\longrightarrow H,
\]
where $\prod_a C_a$ is the finite normalization
of $R/(f)$ and $r$ is the restriction morphism
of local Picard group schemes. These group
schemes are locally of finite type, since
$\dim C_a=N-1\ge3$; as in C.3n,
$G(k)=\operatorname{Pic}(X)$ and
$H(k)=\operatorname{Pic}(Y)$.

For every closed abelian subvariety $B\subset H$
there exist a closed abelian subvariety $A\subset G$,
an integer $e\ge0$, and isogenies of abelian
varieties over $k$
\begin{equation}\label{eq:C3-abelian-local-lift}
 s:B\longrightarrow A,\qquad
 r_A:A\longrightarrow B,\qquad
 r_A\circ s=[p^e]_B,
\end{equation}
where $r_A$ is induced by $r$. Both isogenies
have $p$-primary kernels. Thus every target
abelian subgroup lifts up to a power-of-$p$
isogeny at the level of the specified local
Picard functors.

\begin{proof}
We state separately the two geometric inputs.
First, in the present normal essentially-finite-type
scope, the local Neron--Severi group
\[
                  \operatorname{NS}^{\mathrm{loc}}
                         =G/G^0
\]
has finitely generated point group. This is
the positive-characteristic normal case stated
in \citet[Definition 3.4 and the following
paragraph, p.~468]{D_Kollar2016LocalPicard};
the cited argument combines Boutot's method
with alterations. We only need the weaker
consequence that its prime-to-$p$ torsion
has bounded exponent. Second,
\citet[Theorem 1.7(ii), p.~463, with Definition
3.3]{D_Kollar2016LocalPicard} makes
$\ker(r)$ a unipotent algebraic group.
The target in that theorem is interpreted by
its closed unit section; replacing the divisor
by its finite normalization does not change
the kernel, since its puncture is the same
$Y$. Henselization does not change the local
Picard functor, by its definition using henselian
tensor pairs. If $f$ is only given in $R$, it
descends to a pointed etale neighbourhood of
the original algebraic model. The cited
theorem applies there, giving the same local
functor and restriction map. These inputs are
not asserted here for an arbitrary excellent
ring outside that essentially-finite-type
geometric scope.

Choose a prime $\ell\ne p$ not dividing the
order of the finite prime-to-$p$ torsion
subgroup of $\operatorname{NS}^{\mathrm{loc}}(k)$.
By C.3p every point of $B[\ell^n](k)$ has a
unique preimage in $G(k)[\ell^n]$. Its image
in $\operatorname{NS}^{\mathrm{loc}}(k)$ is zero,
so the preimage belongs to $G^0(k)$. Since
$k$ is perfect, the reduced identity component
\[
                          G_0=(G^0)_{\mathrm{red}}
\]
is a smooth connected commutative algebraic
group with $G_0(k)=G^0(k)$.
The image $I=r(G_0)$ is a closed connected
algebraic subgroup of $H$. It contains all
$\ell$-power torsion points of $B$.
Their union is Zariski dense in $B$: otherwise
its closed subgroup closure would have smaller
identity-component dimension, contradicting
the $\ell^{2n\dim B}$ torsion points in
$B[\ell^n]$ and the corresponding growth in
that smaller-dimensional subgroup. Thus
\[
                              B\subset I.
\]

Let $E$ be the reduced identity component
of $r^{-1}(B)\cap G_0$. Then $E\to B$ is
surjective. Indeed the whole preimage maps
onto $B$, and the images of its finitely
many components are translates of the
closed connected image of $E$; a finite
union of proper closed subvarieties cannot
cover the irreducible variety $B$. A
surjective homomorphism of algebraic groups
is faithfully flat; here this also follows
directly from generic flatness followed
by translations. Write its kernel as $U$.
It is a closed subgroup scheme of the
unipotent group $\ker(r)$, hence is
unipotent and of finite type.

There is an integer $e$ such that
$[p^e]_U=0$ as a morphism of group schemes.
For example, a finite-type unipotent group
over the algebraically closed field embeds
in a group of upper unitriangular matrices;
if the matrices have size $d$ and $p^e\ge d$,
then $(1+M)^{p^e}=1+M^{p^e}=1$ in every
coefficient algebra. This argument includes
nonreduced kernels. The commutativity of $E$
now makes $[p^e]_E$ constant on the $U$-cosets.
Faithfully flat descent through $E\to B=E/U$
therefore gives a unique homomorphism
\[
                         \widetilde s:B\longrightarrow E
\]
with $\widetilde s\circ r_E=[p^e]_E$.
Composing with $r_E$ and using its faithful
flatness yields
\[
                            r_E\circ\widetilde s=[p^e]_B.
\]

Let $A$ be the scheme-theoretic image of
$\widetilde s$ in $E\subset G$. It is
proper, connected and reduced, since $B$
has those properties. A reduced algebraic
group over the perfect field $k$ is smooth,
so $A$ is an abelian variety. The kernel
of $\widetilde s$ is contained in
$B[p^e]$, hence is finite and $p$-primary.
Thus $s:B\to A$ is an isogeny. Its
composition with $r_A$ is $[p^e]_B$, so
$r_A$ is also surjective with finite
$p$-primary kernel, proving
\eqref{eq:C3-abelian-local-lift}.
\end{proof}

The exact ordinary-equivalence consequence
is conditional on an additional family
assertion. If the points of the resulting
$A$ are represented by ordinary algebraic
equivalences on the fixed product $X$ times
curves, then $r_A$ makes the corresponding
target classes lift in ordinary
$\operatorname{Pic}_{\mathrm{alg}}$, and hence
in $J^1$ with rational coefficients.
The theorem itself gives a morphism into
the henselized-product local Picard functor;
that alone is not this additional family
assertion. In particular the proof has not
used a universal actual line bundle on
$X\times A$, a blanket henselian GAGA
essential-surjectivity statement, or an
uncontrolled norm from an etale neighbourhood.
The further ordinary-equivalence step is
the realization of the lifted points by
actual product families. C.3u supplies
a direct torsion argument for smooth-cone
sources in every characteristic. C.3w
settles point realization for the present
algebraizable scope using alterations,
proper curve families, and transfer.
It consequently proves ordinary
surjectivity in the characteristic-$p$
scope of this paragraph, while retaining
the distinction between point realization
and a universal actual family on the
whole local Picard subgroup.
